% answers_short.tex —— 2014 年·北京卷·理科 · 简短答案速查
% 由 tools/make_short_answers.py 自 answers.tex 的 \gksolhead/\gkanswer 提取，
% 并与 2014/2014_数学真题_北京理_答案要点.md 逐题比对后生成，勿手改（改源 MD 后重跑）。

\begin{gknotice}
\textbf{使用说明}：本节只给最终答案便于快速核对，完整推导见解析版。答案为编写者独立推导，\textbf{非官方标准答案}。
\end{gknotice}

\gkpart{选择题}
\begin{center}\begin{tabular}{ccccccccc}
\toprule
题号 & 1 & 2 & 3 & 4 & 5 & 6 & 7 & 8 \\
\midrule
答案 & C & A & B & C & D & D & D & B \\
\bottomrule
\end{tabular}\end{center}

\gkpart{填空题}
\begin{itemize}[leftmargin=2.6em]
  \item[\textbf{9.}] $-1$。$\dfrac{1+\mathrm i}{1-\mathrm i}=\mathrm i$，平方得 $-1$。
  \item[\textbf{10.}] $\sqrt5$。由 $\lambda\boldsymbol a=-\boldsymbol b$ 得 $|\lambda|=|\boldsymbol b|=\sqrt{2^2+1^2}=\sqrt5$。
  \item[\textbf{11.}] $\dfrac{x^2}3-\dfrac{y^2}{12}=1$；$y=\pm2x$。设 $\dfrac{y^2}4-x^2=\lambda$，代入 $(2,2)$ 得 $\lambda=-3$。
  \item[\textbf{12.}] $8$。由条件得 $a_8>0$，$a_9<0$，故前 $8$ 项和最大。
  \item[\textbf{13.}] $36$。$A$，$B$ 相邻共 $2\times4!=48$ 种，其中 $A$ 也与 $C$ 相邻的有 $2\times3!=12$ 种，相减得 $36$。
  \item[\textbf{14.}] $\pi$。由相角条件只能取 $\omega=2$（$\omega=3$ 在题给区间上不单调）。
\end{itemize}

\gkpart{解答题}
\begin{itemize}[leftmargin=2.6em]
  \item[\textbf{15.}] (1) $\dfrac{3\sqrt3}{14}$。$\angle ADC=\dfrac\pi3+\angle BAD$，用两角差的正弦展开；(2) $BD=3$，$AC=7$。正弦定理得 $BD=3$，再由余弦定理得 $AC^2=49$
  \item[\textbf{16.}] (1) $\dfrac12$。$10$ 场中命中率超过 $0.6$ 的有 $5$ 场；(2) $\dfrac{13}{25}$。主场中 $3$ 场、客场中 $2$ 场超过 $0.6$，恰一场超过的配对数 $3\times3+2\times2=13$；(3) $E(X)=\bar x$。各场等可能，故 $E(X)$ 就是 $10$ 个命中次数的平均数 $11.4$
  \item[\textbf{17.}] (1) 见解析版；(2) 所成角为 $\dfrac\pi6$，$PH=2$。以 $A$ 为原点建系，平面 $ABF$ 的法向量为 $(0,-1,1)$，$BC=(1,1,0)$，由 $\sin\theta=\dfrac12$ 得 $\theta=\dfrac\pi6$；$H$ 分 $PC$ 为 $2:1$，故 $PH=\dfrac23|PC|=2$
  \item[\textbf{18.}] (1) 见解析版；(2) $a$ 的最大值为 $\dfrac2\pi$，$b$ 的最小值为 $1$。$\dfrac{\sin x}x$ 在 $\left(0,\dfrac\pi2\right)$ 上递减，值域为 $\left(\dfrac2\pi,1\right)$
  \item[\textbf{19.}] (1) $\dfrac{\sqrt2}2$。$a=2$，$b=\sqrt2$，$c=\sqrt2$；(2) 直线 $AB$ 与圆 $x^2+y^2=2$ 相切。原点到 $AB$ 的距离恒为 $\sqrt2$，等于圆的半径
  \item[\textbf{20.}] (1) $T_1(P)=7$，$T_2(P)=8$。$T_2=1+\max\{7,6\}=8$；(2) $m=a$ 与 $m=d$ 两种情况下均有 $T_2(P)\leqslant T_2(P')$。把两个量都化成 $b+c+\max\{a,d\}$ 与 $a+d+\max\{b,c\}$ 后比较；(3) $P:(4,6)$，$(11,11)$，$(16,11)$，$(11,8)$，$(5,2)$；$T_5(P)=52$。逐项计算得 $10,26,42,50,52$，穷举 $120$ 种排列可验证 $52$ 为最小值
\end{itemize}
